\section{Conclusão}

Este trabalho apresentou um estudo piloto sobre o uso de aprendizado de máquina
para reduzir o número de configurações testadas no VQE. Modelos preditivos estimam
o erro energético de cada candidato e ordenam as opções antes da execução, com
validação cruzada agrupada por geometria. O Random Forest obteve o menor erro
absoluto de regressão, enquanto a MLP atingiu 85,2\% de cobertura no cenário
\textit{top}-3 — o teto máximo deste conjunto — superando em 7,4 pontos percentuais
as estratégias estáticas. Essa vantagem equivale a apenas duas geometrias adicionais,
e os intervalos de \textit{bootstrap} incluem zero, de modo que a superioridade da
MLP permanece estatisticamente incerta. Os testes de estresse mostraram, ainda, que o
modelo depende do nome da molécula: sem essa informação, ou diante de uma molécula
inédita, o desempenho caiu ao nível das estratégias simples, indicando memorização de
padrões por família e não aprendizado de leis gerais da química quântica. A principal
contribuição deste piloto é demonstrar que simulações passadas podem guiar
inteligentemente novos casos correlacionados, eliminando 83,3\% dos candidatos da
busca em grade — embora essa redução não se traduza linearmente em economia de tempo,
dado o custo distinto de cada circuito e o investimento inicial de coleta dos dados de
treino.

Como trabalhos futuros, pretendemos ampliar o estudo com mais moléculas e bases
atômicas, múltiplas sementes de inicialização e modelos de ruído para aproximar as
simulações de hardware NISQ real. Também investigaremos atributos físico-químicos
universais que dispensem o nome categórico da molécula, uma avaliação multiobjetivo
que equilibre erro, custo do circuito e tempo de execução, e o ponto de equilíbrio
financeiro entre investimento no treino e ganho nas previsões subsequentes. O
código-fonte e os dados serão disponibilizados publicamente após a revisão
duplo-anônima.